\begin{thebibliography}{99}
\bibitem{CGT08} Ciabattoni, A., N. Galatos, and K. Terui, From axioms to
  analytic rules in nonclassical logics, in \emph{LICS 2008}, IEEE, 2008,
  pp.~229--240.
\bibitem{GJKO07} Galatos, N., P. Jipsen, T. Kowalski, and H. Ono,
  \emph{Residuated Lattices: An Algebraic Glimpse at Substructural Logics},
  Elsevier, 2007.
\bibitem{Jer16} Jeřábek, E., A note on the substructural hierarchy,
  \emph{Mathematical Logic Quarterly} 62(1--2):102--110, 2016.
  arXiv:1507.00700.
\bibitem{Min90} Mints, G., Gentzen-type systems and resolution rules.
  Part~I: propositional logic, in \emph{COLOG-88}, LNCS 417, Springer, 1990,
  pp.~198--231. \plot{check}
\bibitem{PG86} Plaisted, D.~A., and S. Greenbaum, A structure-preserving
  clause form translation, \emph{Journal of Symbolic Computation}
  2(3):293--304, 1986.
\bibitem{Ryb97} Rybakov, V.~V., \emph{Admissibility of Logical Inference
  Rules}, Elsevier, 1997.
\bibitem{Sve03} Švejdar, V., On the polynomial-space completeness of
  intuitionistic propositional logic, \emph{Archive for Mathematical Logic}
  42(7):711--716, 2003.
\bibitem{Tse68} Tseitin, G.~S., On the complexity of derivation in
  propositional calculus, in \emph{Studies in Constructive Mathematics and
  Mathematical Logic, Part II}, 1968, pp.~115--125.
\bibitem{Dum59} Dummett, M., A propositional calculus with denumerable
  matrix, \emph{The Journal of Symbolic Logic} 24(2):97--106, 1959.
\bibitem{Jan68} Jankov, V.~A., On the extension of the intuitionist
  propositional calculus to the classical calculus, and the minimal calculus
  to the intuitionist calculus, \emph{Mathematics of the USSR-Izvestiya}
  2(1):205--208, 1968.
\bibitem{Kas05} Kashima, R., Problems on axiomatization of intermediate
  propositional logics, in \emph{the 39th MLG meeting}, 2005, pp.~59--62.
\bibitem{Kom05} Komori, Y., A problem on logics axiomatized with formulas
  minimal in classical logic, and more (in Japanese), MSJ Autumn Meeting,
  2005.
\bibitem{NM21} Nakamura, Y., and N. Matsuda, On implicational intermediate
  logics axiomatizable by formulas minimal in classical logic: a
  counter-example to the Komori--Kashima problem, \emph{Studia Logica}
  109:1413--1422, 2021.
\bibitem{Sta79} Statman, R., Intuitionistic propositional logic is
  polynomial-space complete, \emph{Theoretical Computer Science}
  9:67--72, 1979.
\bibitem{Lean} de Moura, L., and S. Ullrich, The Lean 4 theorem prover
  and programming language, in \emph{CADE-28}, LNCS 12699, 2021,
  pp.~625--635.
\end{thebibliography}
